\section{Introduction}
\label{sec:intro}
%
In recent years, diffusion models~\citep{song2020score, ho2020denoising, sohldickstein2015deep, song_improved_2020, song_generative_2020} have achieved state of the art performance across diverse modalities, including image~\citep{dhariwal_diffusion_2021, rombach2022high, esser_scaling_2024}, audio~\citep{popov_grad-tts_2021, jeong_diff-tts_2021, huang_prodiff_2022, lu_conditional_2022}, and video~\citep{ho_imagen_2022, ho_video_2022, blattmann_align_2023, wu_tune--video_2023}.
%
These models belong to a broader class of approaches including flow matching~\citep{lipman2022flow}, rectified flow~\citep{liu2022flow}, and stochastic interpolants~\citep{albergo2022building, albergo2023stochastic}, which construct a path in the space of measures between a base and a target distribution by specifying an explicit mapping between samples from each.
%
This construction yields a dynamical transport equation governing the evolution of the time-dependent probability measure along the path.
%
The generative modeling problem then reduces to learning a velocity field that accomplishes this transport, which provides an efficient and stable training paradigm.

At sample generation time, models in this class generate data by smoothly transforming samples from the base into samples from the target through numerical integration of a differential equation.
%
While effective, the number of integration steps required to produce high-quality samples incurs a cost that can limit real-time applications~\citep{chi_diffusion_2024}.
%
Comparatively, one-step models such as generative adversarial networks~\citep{goodfellow_generative_2014, goodfellow_generative_2020, creswell_generative_2018} are notoriously difficult to train~\citep{metz_unrolled_2017, arjovsky_wasserstein_2017}, but can be orders of magnitude more efficient to sample, because they only require a single network evaluation.
%
As a result, there has been significant recent research effort focused on maintaining the stable training of diffusions while reducing the computational burden of inference~\citep{karras_elucidating_2022}.
%
\begin{figure}[t]
	\centering
	\includegraphics[width=0.85\textwidth]{figs/overview/fmm_t_labeled_less_compressed.pdf}
	\caption{
		\textbf{Overview of Flow Map Matching.}
		%
		Our approach learns the two-time flow map $X_{s, t}$ that transports the solution of an ordinary differential equation from time $s$ to time $t$.
		%
		Unlike methods that learn instantaneous velocity fields, this bidirectional map can be used to build an integrator with arbitrary discretization.
		%
		The integrator is exact in theory, and --crucially-- its number of steps can be adjusted post-training to balance accuracy and computational efficiency.
		The flow map can be distilled from a known velocity field or learned directly, and supports arbitrary base distributions, as illustrated here with image-to-image translation.
	}
	\label{fig:overview}
\end{figure}
%
Towards this goal, we introduce Flow Map Matching---a theoretical framework for learning the two-time flow map of a probability flow.
%
Specifically, our primary contributions can be summarized as:
%

\begin{itemize}[leftmargin=0.2in]
	\item \textbf{A mathematical framework.}
	      %
	      We characterize the flow map theoretically by identifying its fundamental mathematical properties, which we show immediately lead to distillation-based algorithms for estimating the flow map from a pre-trained velocity field.
	      %
	      Using our characterization, we establish connections between our framework and recent approaches for learning consistency models~\citep{song_consistency_2023, song_improved_2023, kim_consistency_2024} and for progressive distillation of a diffusion model~\citep{salimans_progressive_2022, zheng_fast_2023}.
	      %
	\item \textbf{Lagrangian vs Eulerian learning.}
	      %
	      % We develop a novel \textit{Lagrangian} objective for distilling a flow map from a pre-trained velocity field, and show that it outperforms a related \textit{Eulerian} loss, which we prove is the continuous-time limit of consistency distillation~\citep{song_consistency_2023}. We show that both the Lagrangian and Eulerian losses control the Wasserstein distance between the teacher and the student models.
	      We develop a novel \textit{Lagrangian} objective for distilling a flow map from a pre-trained velocity field, which performs well across our experiments. This loss is complementary to a related \textit{Eulerian} loss, which we prove is the continuous-time limit of consistency distillation~\citep{song_consistency_2023}. We show that both the Lagrangian and Eulerian losses control the Wasserstein distance between the teacher and the student models.
	      %
	\item \textbf{Direct training.}
	      %
	      We extend our Lagrangian loss to design a novel direct training objective for flow maps that eliminates the need for a pre-trained velocity field.
	      %
	      We also discuss how to modify the Eulerian loss to allow for direct training of the map, and identify some caveats with this procedure.
	      %
	      Our approach elucidates how the standard training signal for flow matching can be adapted to capture the more global signal needed to train the flow map.
	      %
	\item \textbf{Numerical experiments.}
	      %
	      We validate our theoretical framework through experiments on CIFAR-10 and ImageNet-$32$, where we achieve high sample quality at significantly reduced computational cost.
	      %
	      In particular, we highlight that learning the two-time flow map enables post-training adjustment of sampling steps, allowing practitioners to systematically trade accuracy for computational efficiency.
	      %
\end{itemize}
